\chapter{DATA COLLECTION FRAMEWORK}

\section{System Architecture}

The data collection system follows a \textbf{cloud-native, serverless architecture}
designed to minimize operational overhead while ensuring daily data availability.
Figure~\ref{fig:architecture} illustrates the overall system design.

\begin{figure}[H]
  \centering
  % \includegraphics[width=0.9\textwidth]{figures/architecture.png}
  \fbox{\parbox{0.85\textwidth}{\centering [Figure: System Architecture Diagram]\\ GitHub Actions → Apify Actors → Python Parsers → Supabase PostgreSQL + Storage}}
  \caption{System architecture overview}
  \label{fig:architecture}
\end{figure}

The pipeline consists of five layers:
\begin{enumerate}
  \item \textbf{Orchestration Layer:} GitHub Actions cron job (06:00 UTC daily) triggers the ingest
  \& analytics pipeline. A second weekly cron (07:00 UTC Monday) runs the
  reviews + sentiment pipeline.
  \item \textbf{Collection Layer:} Apify cloud actors perform web scraping with managed proxies
  \item \textbf{Processing Layer:} Python scripts parse, validate, and enrich raw data
  \item \textbf{Storage Layer:} Supabase PostgreSQL stores structured data; Supabase Storage holds product images
  \item \textbf{Delivery Layer:} A Streamlit dashboard exposes the warehouse
  for ad-hoc exploration. In parallel, the self-hosted OpenClaw runtime
  exposes thirteen read-only Python skills to three specialist Telegram bots,
  delivering both interactive replies and three scheduled cron briefs
  (daily Strategist, weekly Spy, weekly Detective). The delivery layer is
  described in detail in Chapter~4.
\end{enumerate}

\section{Data Sources and Collection Actors}

\subsection{Amazon Category Rankings}
Category ranking data is collected using the \texttt{junglee/Amazon-crawler} actor
(actor ID: \texttt{BG3WDrGdteHgZgbPK}) with the following configuration:

\begin{lstlisting}[language=Python, caption=Category scraper configuration]
{
    "categoryOrProductUrls": [{"url": category_search_url}],
    "maxItemsPerStartUrl": 50,
    "scrapeProductDetails": False,
    "proxyCountry": "AUTO_SELECT_PROXY_COUNTRY"
}
\end{lstlisting}

Setting \texttt{scrapeProductDetails: False} ensures the actor operates in list mode,
efficiently extracting rank, title, price, brand, and sponsored status for the top 50
results without visiting individual product pages.

\subsection{Product Detail Snapshots}
Daily product snapshots for the 30 watchlist ASINs are collected using the same actor
with a search URL workaround: rather than supplying direct \texttt{/dp/ASIN} URLs
(which produced empty datasets due to an actor-level behavior with the default dataset),
ASIN-specific search queries (\texttt{s?k=ASIN}) are used. The actor identifies the
target product as the top search result and navigates to its detail page, correctly
writing product data to the output dataset.

ASINs are processed in batches of 10 to avoid rate limiting, resulting in 3 actor
runs per daily collection cycle.

\subsection{Customer Reviews}
Customer reviews are collected using the \texttt{web\_wanderer/amazon-reviews-extractor}
actor (\$1.00 per 1,000 reviews). Reviews are collected once per week for each
watchlist ASIN, with a maximum of 50 reviews per product sorted by recency, filtered
to English-language reviews only. This actor additionally provides aspect-level
sentiment annotations and an AI-generated review summary per product, which are
stored alongside raw review text.

\section{Database Schema}

All structured data is stored in a Supabase PostgreSQL instance. Table~\ref{tab:schema}
summarizes the core tables.

\begin{longtable}{lll}
  \toprule
  \textbf{Table} & \textbf{Granularity} & \textbf{Key Fields} \\
  \midrule
  \endhead
  \texttt{asins} & Per ASIN & asin, product\_name, brand, category \\
  \texttt{category\_rankings} & Per ASIN per day & rank, price, is\_sponsored, stars \\
  \texttt{daily\_snapshots} & Per ASIN per day & bsr, price, buy\_box\_winner, in\_stock \\
  \texttt{reviews\_raw} & Per review & review\_text, rating, sentiment\_label \\
  \texttt{entrant\_exit\_events} & Per event & event\_type, rank\_today, is\_top10 \\
  \texttt{price\_tier\_daily} & Per ASIN per day & cluster\_label, cluster\_name \\
  \texttt{listing\_quality\_score\_daily} & Per ASIN per day & lqs\_total, sub-scores \\
  \texttt{brand\_momentum\_daily} & Per ASIN per day & bms\_score, bsr\_velocity \\
  \texttt{review\_sentiment\_daily} & Per ASIN per day & avg\_sentiment\_score, ratios \\
  \texttt{image\_change\_events} & Per event & hash\_before, hash\_after, distance \\
  \texttt{content\_change\_events} & Per event & change\_area, change\_level \\
  \texttt{alerts} & Per alert & alert\_type, severity, message \\
  \bottomrule
  \caption{Database schema summary}
  \label{tab:schema}
\end{longtable}

\section{Target Categories and Watchlist}

\subsection{Category Selection}
Three consumer electronics sub-categories were selected based on:
\begin{itemize}
  \item \textbf{Competitive intensity:} High number of competing ASINs and frequent ranking changes
  \item \textbf{Price diversity:} Sufficient price range to support meaningful tiering
  \item \textbf{Data availability:} Consistent BSR and review data across category members
\end{itemize}

Table~\ref{tab:categories} summarizes the selected categories.

\begin{table}[H]
  \centering
  \begin{tabular}{llc}
    \toprule
    \textbf{Category} & \textbf{Search URL} & \textbf{ASINs Tracked} \\
    \midrule
    Gaming Keyboards   & amazon.com/s?k=gaming+keyboard      & 50 + 10 watchlist \\
    True Wireless Earbuds & amazon.com/s?k=true+wireless+earbuds & 50 + 10 watchlist \\
    Portable Chargers  & amazon.com/s?k=portable+charger     & 50 + 10 watchlist \\
    \bottomrule
  \end{tabular}
  \caption{Target categories and collection scope}
  \label{tab:categories}
\end{table}

\subsection{Watchlist ASIN Selection}
Ten ASINs per category were selected as the primary watchlist based on rank
diversity (spanning ranks 1–50), price tier diversity, and brand diversity.
The watchlist includes products from major brands (e.g., Logitech, Razer, Sony,
Anker) as well as emerging direct-to-consumer brands, enabling comparative analysis
across competitive tiers.

\section{Image Processing Pipeline}

Product images are processed through a perceptual hashing pipeline to detect
visual changes:

\begin{enumerate}
  \item The main product image URL is extracted from the actor output
  (prioritizing \texttt{highResolutionImages} over \texttt{thumbnailImage})
  \item The image is downloaded and converted to RGB format using Pillow
  \item A 64-bit pHash is computed using the \texttt{imagehash} library
  \item The image is uploaded to Supabase Storage bucket \texttt{product-images}
  at path \texttt{YYYY-MM-DD/ASIN/main.jpg}
  \item The Hamming distance between today's and yesterday's pHash is computed;
  distances exceeding the threshold of 10 bits trigger an \texttt{image\_changed} flag
\end{enumerate}

\section{Data Quality and Limitations}

Several data quality considerations apply to this study:

\textbf{ASIN resolution:} The search-based collection approach occasionally returns
a related but non-target ASIN as the top search result. Such records are filtered
by validating the returned ASIN against the watchlist before database insertion.

\textbf{BSR availability:} Not all products maintain a BSR at any given time; newly
launched products or those with insufficient recent sales history may return null BSR
values. These are stored as \texttt{NULL} and excluded from velocity calculations.

\textbf{Actor rate limiting:} Apify actor runs are capped at 10 ASINs per run to
avoid triggering Amazon's rate limiting mechanisms, which can result in empty
datasets for larger batch sizes.

\textbf{Collection window:} The 13-day data collection window
(2026-04-17 to 2026-04-29) limits the statistical power of time-series
analyses. Findings should be treated as illustrative of the methodology
rather than definitive market conclusions.
